\documentclass[a4paper,11pt]{article}

\usepackage{revy}
\usepackage[utf8]{inputenc}
\usepackage[T1]{fontenc}
\usepackage[danish]{babel}
\usepackage{amsmath,amssymb,amsfonts,amsthm,textcomp}


\revyname{MatematikRevyen}
\revyyear{2026}
\version{0.1}
\eta{$3$ minutter}
\status{Færdig}

\title{Instruktorshopping}
\author{Leif '21 \& Carl '21}

\begin{document}
\maketitle

\begin{roles}
\role{S1}[Skuespiller] Studerende
\role{S2}[Skuespiller] Studerende
\role{B1}[Skuespiller] Bod-ejer (\textit{meget} sælger)
\role{B2}[Skuespiller] Bod-ejer (almindelig sælger)
\end{roles}

\begin{props}
\item 
\end{props}


\begin{sketch}
\scene{S2 står på trappen, S2 kommer hen til S1}

\says{S1} Hey, skal du med ud og handle ind til den nye blok? 

\says{S2} Uhh, ja. Det har jeg ikke fået gjort endnu. Vi ku tage i svalehallerne? De har et ret godt udvalg.

\says{S1} Ja! Næste blok skulle være hård, så det ville være nice at finde den helt rigtige instruktor. 

\scene{S1 og S2 går op på scenen som er en markedsplads. Der er to boder, hver med 2-4 instruktorer foran. Rigtige statister eller billeder/tegninger}

\says{B1} Kom og køb, kom og køb! KU's bedste instruktorer til en fair pris!

\says{B2} Se herover, se herover! Instruktorer, der er hjælpsomme OG ambitiøse!

\says{S2} Det er lidt overvældende med alle de muligheder

\says{S1} Lad os bare starte et sted

\scene{S1 og S2 går over mod B1}

\says{B1} Veeelkommen! Hvad er I så på udkig efter i dag?

\says{S2} Jeg leder efter en instruktor til KomAn

\says{B1} KomAn, KomAn, KomIgen, Kom Her-hen! Jeg har \textit{lige} det du har brug for. 

\says{B1}[peger på en af instruktorerne] Ham/hende her er et lille geni. De når altid alle opgaverne OG.... nogle gange også løsningerne til dem. Hahaaa, den var god hva? 

\says{S2} Øh, joh. Men altså, er de gode til at formidle og sådan?

\says{B1} Om de er! Hvis du når igennem hele kurset uden at kunne sige holomorf, så får du alle pengene igen. Ga-ran-te-ret!

\says{S2} Altså, jeg ville nok gerne lære lidt mere end det 

\says{S1} Har du nogen til StatMet?

\says{B1} StatMet, StatGæt, Jørgen Leth, Tag og sæt! dig her. Du ka tro, at jeg har det, men øhh

\says{B1}[læner sig ind mod S1] mellem os... så er der mere vaLue for moNEY med et lille chatgpt+ abonnement til det kursus der

\says{B1}[peger på en af instruktorerne] Det er alligevel det ham/hende her gør

\says{S1} Tak øh, jeg kigger nok bare lidt videre

\scene{S1 går over mod B2, mens S2 og B1 snakker videre (ikke højt)}

\says{B2} Hej du! Hvad skulle det være?

\says{S1} Har du nogle instruktorer til StatMet? Hvad med ham/hende der? (peger på en af dem)

\says{B2} Glimrende valg! Det er tredje gang de er TA i det kursus, og jeg har hørt, at de er begyndt at tage kage med til eftermiddagsmoduler 

\says{B1}[fra den anden side af scene] Er det KAge! du er efter? Jeg har ham/hende her. Skide dårlig til matematik men dygtig bager! Du får halv pris på grund af det faglige! 

\says{B2}[til S1] Du skal ikke tage dig af B1, de siger hvad som helst

\says{B1} Hvad som helst, bla blom blelst! Ha! B2 giver dig en afdanket instruktor, der er TA for tredje gang for at tjene lette penge. Nej! Hem/hende her tager kaffe med til morgenmoduler!

\says{S2}[bryder ind] Uhh, den tager jeg!!

\says{B1} Solgt!

\says{S1}[til S2] Duuude! Det var ikke engang til dit kursus?!?

\says{S2}[forvirret] Var det ik? Er eftertilmeldingen stadig åben?

\scene{S1 ryster på hovedet og vender tilbage til B2. S2 tager en kop kaffe fra instruktoren de lige har købt}

\says{S1}[lidt opgivende] Kan jeg bare få ham/hende der?

\says{B2} Uhhh, det ser ud til at der lige er blevet fyldt, desværre. 

\says{B1}[drille-syngende] Ved du hvem er aaaldrig har fyyyldt???

\scene{S1 vender sig irriteret mod B1 som om de skulle til at sige 'hold kæft' men siger ikke noget og vender sig tilbage mod B2. S2 nyder stadig sin kaffe.}

\says{S1}[opgivende] Hvad med ham/hende der?

\says{B2}[tøvende] Tjo, der er altid plads hos dem. Meeen de forventer, at du har kigget på opgaverne hjemmefra.

\says{S1} Aargh, ved du hvad. Måske jeg bare går på halv tid i stedet.

\scene{Lys ned}

\end{sketch}
\end{document}
